\chapter{Introduction}
\label{ch:introduction}

\section{What fBlockKit is}

\fbk{} is a local toolbox for strongly correlated systems and
lanthanide/actinide calculations. It packages the expert knowledge around two
recurring jobs into deterministic procedures:

\begin{enumerate}
  \item \textbf{Generating input}: you describe the system (elements, charge,
    multiplicity, f-block metal valence, targets, structure), and \fbk{}
    proposes the method chain, the basis set and ECP, a starting active space
    and the convergence settings, then writes a directly runnable ORCA input
    file together with plain run guidance.
  \item \textbf{Reading results}: you give it an output file (or a structure
    file), and \fbk{} reports what the calculation actually did: orbital
    composition, occupation and entropy analysis, multi-reference character,
    local spin distribution, crystal-field parameters, cross-level solution
    consistency, SCF problems, transition-state checks -- each conclusion with
    its provenance and, where needed, a refusal to conclude.
\end{enumerate}

The program is a workbench, not an engine. It wraps knowledge around existing
codes (ORCA is the deep target of version 0.1; an OpenMolcas magnetic-property
chain template is included), and it is designed for chemists who do not
program: a numbered menu, prompts that always show their default, and a
``record once, replay forever'' script model.

\section{What fBlockKit is not}

\begin{refusalbox}
\begin{itemize}
  \item Its core does \textbf{not} run calculations. You submit the generated
    input to ORCA (or OpenMolcas) yourself -- except that the optional execution
    bridge, off by default and acting only when explicitly asked, can run one
    generated input locally (menu 50); the core never queues a job on any
    scheduler and never waits for one.
  \item It does \textbf{not} use machine learning in its decision layer. Every
    recommendation and every diagnosis is a deterministic rule with declared
    conditions, an action and provenance; the same input always produces the
    same output.
  \item It does \textbf{not} use the network. Nothing is uploaded, nothing is
    fetched; the program works in an offline environment.
  \item It does \textbf{not} modify your files. Every product is a new file
    beside the input (\file{*.fbk.md}, \file{*.fbk.json}, \file{*.fbk.inp},
    \file{*.fix\_*.inp}); the originals are only ever read.
\end{itemize}
\end{refusalbox}

\section{Design principles}

\begin{description}[leftmargin=1.6em]
  \item[Refuse rather than guess.] When the data needed for a conclusion is
    missing, or a request is structurally impossible (a 5f-in-core
    pseudopotential for an f--f transition, for example), the program says so
    and states the next step. A suspicious number is worse than no number.
  \item[Everything carries provenance.] Three classes are marked in every
    report: a program manual (§ number and URL), a literature reference (full
    citation with DOI, plus a paste-ready BibTeX entry), or a measured record
    of this project (a file path or a fixture). Literature entries are checked
    against a bundled BibTeX database at load time, so a rule whose source
    cannot be resolved does not load at all.
  \item[The interaction is a script.] Every sequence of answers can be saved
    (or recorded while you type) and replayed: the same script produces the
    same output, byte for byte. This is also the program's regression
    acceptance criterion.
  \item[Registration is not use.] The built-in index of external tools records
    two independent facts for every entry: how \fbk{} intends to relate to the
    tool (absorb / interface / index) and whether the tool is actually used
    \emph{now} (\code{active} / \code{registered}). Do not read the index as a
    feature list.
\end{description}

\section{Scope of version 0.1}

Version 0.1 covers ORCA deeply (input generation and output parsing) and
OpenMolcas only through the magnetic-property chain template. The analysis and
diagnosis side reads ORCA 6.x output; the parser patterns are fixed against
registered real output files of ORCA 6.1.1 (\file{fixtures/orca/}). Other
programs (Gaussian, PySCF, \dots) are on the long-term roadmap, not in this
version.

The program is licensed under Apache-2.0 and is in early development: the
report formats and interfaces may still change between releases. When you use
it, cite it as described in \file{CITATION.cff}.
